% Chapter 6. Source pages PDF 149--178 (printed 141--170).
\section{Limits via representables and adjoints}
\begin{proposition}[Limits represented by cones]
\label{prop:lbc-limit-representable}
\source{leinster-basic-category-theory:page-0150}
\dcref{6.1.1}
For a small $I$, a category $\mathcal A$, and $D:I\to\mathcal A$, limit cones
on $D$ correspond bijectively to representations of
$\operatorname{Cone}(-,D):\mathcal A^{\op}\to\Set$; the representing objects are
the vertices of the limit cones.
\uses{def:lbc-limit,def:lbc-representable-cov}
\end{proposition}
\begin{proof}
The displayed map sends a mediating arrow to its induced cone. Bijectivity is
exactly existence and uniqueness in the definition of a limit.
\end{proof}

\begin{corollary}[Uniqueness of limits]
\label{cor:lbc-limit-unique}
\source{leinster-basic-category-theory:page-0151}
\dcref{6.1.2}
Limits of a fixed diagram are uniquely isomorphic.
\uses{prop:lbc-limit-representable,def:lbc-isomorphism}
\end{corollary}
\begin{proof}
Apply the bijection in Proposition~\ref{prop:lbc-limit-representable} to the
two limiting cones to obtain unique arrows each way; uniqueness forces their
composites to be identities.
\end{proof}

\begin{lemma}[Natural transformations induce cone maps]
\label{lem:lbc-natural-cone-map}
\source{leinster-basic-category-theory:page-0151}
\dcref{6.1.3}
For a natural transformation $\alpha:D\Rightarrow D'$ and a limit cone on $D'$, a
limit cone on $D$ is obtained by composing its legs with $\alpha$; the induced
map between limit objects is unique and natural.
\uses{def:lbc-limit,def:lbc-natural-transformation}
\end{lemma}
\begin{proof}
Compose each leg with the corresponding component of $\alpha$; naturality of
$\alpha$ gives a cone. The universal property of the target limit gives the
unique mediating map, and uniqueness proves compatibility with further maps.
\end{proof}

\begin{proposition}[Limits as right adjoints]
\label{prop:lbc-limit-right-adjoint}
\source{leinster-basic-category-theory:page-0152}
\dcref{6.1.4}
If $\mathcal A$ has all limits of shape $I$, the diagonal functor
$\Delta:\mathcal A\to[ I,\mathcal A]$ has a right adjoint $\limx$ sending a diagram
to its limit. Conversely, existence of this right adjoint gives all such limits.
\uses{def:lbc-limit,def:lbc-adjunction,con:lbc-functor-category}
\end{proposition}
\begin{proof}
The adjunction bijection is the limit universal property
$[I,\mathcal A](\Delta A,D)\cong\mathcal A(A,\limx D)$. The converse reads the
unit/counit of the adjunction as a universal cone.
\end{proof}

\section{Limits and colimits of presheaves}
\begin{lemma}[Cones as set-theoretic limits]
\label{lem:lbc-pointwise-set-limit}
\source{leinster-basic-category-theory:page-0154}
\dcref{6.2.1}
Let $I$ be small, $\mathcal A$ locally small, $D:I\to\mathcal A$ a diagram,
and $A\in\mathcal A$. Then
\[\operatorname{Cone}(A,D)\cong\limx_{i\in I}\mathcal A(A,D(i)),\]
naturally in $A$ and $D$.
\uses{def:lbc-category,def:lbc-limit,def:lbc-family-product}
\end{lemma}
\begin{proof}
The set-theoretic limit consists of families $(f_i)$ satisfying
$D(u)f_i=f_j$ for every $u:i\to j$, which is exactly the cone condition.
\end{proof}

\begin{proposition}[Representables preserve limits]
\label{prop:lbc-representables-preserve-limits}
\source{leinster-basic-category-theory:page-0155}
\dcref{6.2.2}
For a locally small category, each representable functor preserves all limits
that exist in its domain.
\uses{def:lbc-representable-cov,def:lbc-limit}
\end{proposition}
\begin{proof}
Maps into a limit object correspond by its universal property to compatible
families of maps into the diagram objects, which is exactly the limit of the
representable set-valued diagram.
\end{proof}

\begin{remark}[Representable limit interpretation]
\label{rem:lbc-representable-limit}
\source{leinster-basic-category-theory:page-0155}
\dcref{6.2.3}
Proposition~\ref{prop:lbc-representables-preserve-limits} says that each
representable turns a limit cone into the corresponding set-theoretic limit.
\uses{prop:lbc-representables-preserve-limits}
\end{remark}

\begin{notation}[Evaluation functors]
\label{not:lbc-evaluation}
\source{leinster-basic-category-theory:page-0156}
\dcref{6.2.4}
For categories $\mathcal A,\mathcal S$ and $A\in\mathcal A$, evaluation is the
functor $\operatorname{ev}_A:[\mathcal A,\mathcal S]\to\mathcal S$,
$D\mapsto D(A)$.
\uses{con:lbc-functor-category}
\end{notation}

\begin{theorem}[Limits in functor categories]
\label{thm:lbc-functor-category-limits}
\source{leinster-basic-category-theory:page-0156}
\dcref{6.2.5}
If $\mathcal A$ and $I$ are small and $\mathcal S$ has limits of shape $I$, then
$[\mathcal A,\mathcal S]$ has limits of shape $I$, computed pointwise.
\uses{lem:lbc-pointwise-set-limit,def:lbc-limit,con:lbc-functor-category}
\end{theorem}
\begin{proof}
Evaluate a cone at each $A\in\mathcal A$ and take its limit in $\mathcal S$.
Functoriality follows from uniqueness of mediating maps; the pointwise
universal maps assemble to the universal map in the functor category.
\end{proof}

\begin{corollary}[Presheaf categories are complete]
\label{cor:lbc-presheaf-complete}
\source{leinster-basic-category-theory:page-0158}
\dcref{6.2.6}
Let $I$ and $\mathcal A$ be small and let $\mathcal S$ be locally small. If
$\mathcal S$ has all limits (respectively colimits) of shape $I$, then so does
$[\mathcal A,\mathcal S]$, and each evaluation functor preserves them.
\uses{thm:lbc-functor-category-limits}
\end{corollary}

\begin{warning}[Pointwise limits require hypotheses]
\label{warn:lbc-pointwise-hypothesis}
\source{leinster-basic-category-theory:page-0158}
\dcref{6.2.7}
If the target category lacks limits of shape $I$, the functor category need not
have them.
\uses{thm:lbc-functor-category-limits}
\end{warning}

\begin{proposition}[Limits commute with limits]
\label{prop:lbc-limits-commute}
\source{leinster-basic-category-theory:page-0158}
\dcref{6.2.8}
Let $I,J$ be small and let $\mathcal S$ be locally small with limits of shapes
$I$ and $J$. For every $D:I\times J\to\mathcal S$,
\[\limx_J\limx_I D\cong\limx_{I\times J}D\cong\limx_I\limx_J D,\]
and all these limits exist; in particular $\mathcal S$ has limits of shape
$I\times J$.
\uses{def:lbc-limit,thm:lbc-functor-category-limits}
\end{proposition}
\begin{proof}
Both iterated limits represent the same set of compatible doubly indexed
families, namely the limit over $I\times J$; uniqueness of representing objects
gives the canonical isomorphism.
\end{proof}

\begin{warning}[Dual statement]
\label{warn:lbc-colimits-commute}
\source{leinster-basic-category-theory:page-0160}
\dcref{6.2.10}
The dual of Proposition~\ref{prop:lbc-limits-commute} states that colimits
commute with colimits when the required colimits exist.
\uses{prop:lbc-limits-commute,def:lbc-colimit}
\end{warning}

\begin{corollary}[Presheaf limits]
\label{cor:lbc-presheaf-all-limits}
\source{leinster-basic-category-theory:page-0160}
\dcref{6.2.11}
For a small category $\mathcal A$, the presheaf category
$[\mathcal A^{\op},\Set]$ has all limits and colimits, and each evaluation
functor preserves them.
\uses{cor:lbc-presheaf-complete}
\end{corollary}

\begin{corollary}[Yoneda preserves limits]
\label{cor:lbc-yoneda-preserves-limits}
\source{leinster-basic-category-theory:page-0160}
\dcref{6.2.12}
The Yoneda embedding preserves all limits that exist in a small category.
\uses{cor:lbc-presheaf-complete,prop:lbc-representables-preserve-limits,def:lbc-yoneda-embedding}
\end{corollary}

\begin{definition}[Category of elements]
\label{def:lbc-category-elements}
\source{leinster-basic-category-theory:page-0162}
\dcref{6.2.16}
For a presheaf $X$ on $\mathcal A$, its category of elements $\int X$ has
objects $(A,x\in X(A))$ and maps $(A,x)\to(B,y)$ given by $f:A\to B$ with
$X(f)(y)=x$.
\uses{def:lbc-presheaf}
\end{definition}

\begin{warning}[Yoneda does not preserve colimits]
\label{warn:lbc-yoneda-colimits}
\source{leinster-basic-category-theory:page-0161}
\dcref{6.2.14}
Although Yoneda preserves limits, it need not preserve colimits.
\uses{def:lbc-yoneda-embedding,def:lbc-colimit}
\end{warning}

\begin{theorem}[Density]
\label{thm:lbc-density}
\source{leinster-basic-category-theory:page-0162}
\dcref{6.2.17}
For small $\mathcal A$ and presheaf $X$, the canonical cocone from representables
$H_A$ indexed by $\int X$ has colimit $X$.
\uses{def:lbc-category-elements,def:lbc-representable-contra,def:lbc-colimit}
\end{theorem}
\begin{proof}
At $B$, the coproduct of pairs $(A,x,f:B\to A)$ modulo the relation induced by
maps in $\int X$ identifies with $X(B)$ by $(A,x,f)\mapsto X(f)(x)$. The
relations give exactly equality in $X(B)$, and the identification is natural.
\end{proof}

\begin{remark}[Category of elements and density]
\label{rem:lbc-category-elements-density}
\source{leinster-basic-category-theory:page-0164}
\dcref{6.2.19}
Objects of $\int X$ are generalized elements of the presheaf $X$ of
representable shape. The density theorem says that representables generate
every presheaf by the displayed colimit; this is the categorical meaning of
the term ``dense''.
\uses{def:lbc-category-elements,def:lbc-generalized-element,thm:lbc-density}
\end{remark}

\section{Adjoints and limits}
\begin{theorem}[Left adjoints preserve colimits]
\label{thm:lbc-left-adjoint-colimits}
\source{leinster-basic-category-theory:page-0166}
\dcref{6.3.1}
Every left adjoint preserves all colimits; every right adjoint preserves all
limits.
\uses{def:lbc-adjunction,def:lbc-colimit,def:lbc-limit}
\end{theorem}
\begin{proof}
Apply the adjunction hom-set bijection to the universal cocone: maps from
$F(C)$ correspond to maps from $C$, and the latter are uniquely determined by
their restrictions to the diagram. The dual argument proves preservation of
limits by right adjoints.
\end{proof}

\begin{definition}[Complete category]
\label{def:lbc-complete}
\source{leinster-basic-category-theory:page-0167}
\dcref{6.3.6}
A category is complete when it has all small limits.
\uses{def:lbc-limit}
\end{definition}

\begin{proposition}[Adjoint functor theorem for ordered sets]
\label{prop:lbc-ordered-aft}
\source{leinster-basic-category-theory:page-0168}
\dcref{6.3.7}
For an order-preserving map $G:B\to A$, a left adjoint exists exactly when for
each $a\in A$ the set $\{b\mid a\le G(b)\}$ has a least element.
\uses{def:lbc-adjunction}
\end{proposition}
\begin{proof}
The adjunction condition $F(a)\le b\iff a\le G(b)$ says precisely that $F(a)$
is least among those $b$ satisfying $a\le G(b)$.
\end{proof}

\begin{definition}[Weakly initial set]
\label{def:lbc-weakly-initial-set}
\source{leinster-basic-category-theory:page-0170}
\dcref{6.3.9}
A set $S$ of objects of $\mathcal C$ is weakly initial if every object of
$\mathcal C$ receives at least one map from some $S\in S$.
\uses{def:lbc-category}
\end{definition}

\begin{theorem}[General adjoint functor theorem]
\label{thm:lbc-gaft}
\source{leinster-basic-category-theory:page-0170}
\dcref{6.3.10}
Let $G:\mathcal B\to\mathcal A$ with $\mathcal B$ complete and locally small.
If each comma category $(A\downarrow G)$ has a small weakly initial set, then
$G$ has a left adjoint if and only if $G$ preserves limits.
\uses{def:lbc-weakly-initial-set,def:lbc-complete,def:lbc-comma,cor:lbc-left-adjoint-criterion}
\end{theorem}
\begin{proof}
The forward implication is Theorem~\ref{thm:lbc-left-adjoint-colimits}. For the
converse, preservation of limits makes each comma category complete. In each,
take the limit of the full subcategory on a weakly initial set; the appendix
shows this limit is initial. The comma-object criterion then yields a left
adjoint.
\end{proof}

\begin{theorem}[Special adjoint functor theorem]
\label{thm:lbc-saft}
\source{leinster-basic-category-theory:page-0171}
\dcref{6.3.13}
Let $\mathcal A,\mathcal B$ be locally small, let $\mathcal B$ be complete and
satisfy the standard SAFT hypotheses, and let $G:\mathcal B\to\mathcal A$.
Then $G$ has a left adjoint if and only if $G$ preserves limits.
\uses{thm:lbc-gaft,def:lbc-complete,def:lbc-preserves-limits}
\notready
\end{theorem}
\begin{proof}
\notready
\end{proof}

\begin{definition}[Cartesian closed category]
\label{def:lbc-cartesian-closed}
\source{leinster-basic-category-theory:page-0172}
\dcref{6.3.15}
A category is cartesian closed if it has finite products and, for every $B$, the
functor $(-)\times B$ has a right adjoint, written $(-)^B$.
\uses{def:lbc-product,def:lbc-adjunction}
\end{definition}

\begin{remark}[Internal hom in vector spaces]
\label{rem:lbc-vect-not-closed}
\source{leinster-basic-category-theory:page-0173}
\dcref{6.3.19}
For vector spaces $V,W$, the hom-set is a vector space, but tensoring with a
fixed vector space has no right adjoint in $\Vect_k$ in general; thus $\Vect_k$
is not cartesian closed.
\uses{def:lbc-cartesian-closed}
\end{remark}

\begin{theorem}[Presheaf categories are cartesian closed]
\label{thm:lbc-presheaf-cartesian-closed}
\source{leinster-basic-category-theory:page-0174}
\dcref{6.3.20}
For every small category $\mathcal A$, the presheaf category
$[\mathcal A^{\op},\Set]$ is cartesian closed.
\uses{cor:lbc-presheaf-complete,def:lbc-cartesian-closed}
\end{theorem}
\begin{proof}
Products are pointwise. For presheaves $X,Y$, define $Y^X(A)$ to be the set of
natural transformations $H_A\times X\to Y$. Yoneda and pointwise products
give natural bijections $\operatorname{Nat}(Z\times X,Y)\cong
\operatorname{Nat}(Z,Y^X)$, so $Y^X$ is the exponential.
\end{proof}
